% !TITLE 哀江南赋（节选）
% !AUTHOR 庾信
% !DYNASTY 南北朝
% !GENRE 骈赋
% !ORDER 4

\begin{poem}
日暮途远\textsuperscript{1}，人间何世！
将军一去，大树飘零\textsuperscript{2}；
壮士不还，寒风萧瑟\textsuperscript{3}。

舟楫路穷\textsuperscript{4}，星汉\textsuperscript{5}非乘槎\textsuperscript{6}可上；
风飙道阻，蓬莱\textsuperscript{7}无可到之期。

提挈\textsuperscript{8}老幼，关河\textsuperscript{9}累年。
死生契阔\textsuperscript{10}，不可问天。
况复零落将尽，灵光岿然\textsuperscript{11}！
\end{poem}

\begin{annotations}
\notetext{1}{日暮途远：太阳快落山了，前方的路还很遥远。典出《史记》，伍子胥说自己“日暮途远，故倒行而逆施之”。这里用来自况晚年被困异乡、归途漫漫。}
\notetext{2}{将军一去，大树飘零：指东汉的“大树将军”冯异。他为人谦虚，诸将争功时他总是独自坐在大树下，故得此名。这里比喻优秀的将领们已经离世，如同大树凋零。}
\notetext{3}{壮士不还，寒风萧瑟：化用荆轲刺秦的典故。荆轲渡易水时唱“风萧萧兮易水寒，壮士一去兮不复还”。这里悼念那些一去不返的战死者。}
\notetext{4}{舟楫路穷：船和桨都到了路的尽头，再也无法前进。比喻走投无路。}
\notetext{5}{星汉：银河、天河。天上的星河并非凡人能及的地方。}
\notetext{6}{槎：木筏。传说中有人乘木筏顺海而上，竟然到了天河。庾信说“非乘槎可上”——但不是谁都能乘筏上天河，意味着归乡之路不可能实现。}
\notetext{7}{蓬莱：传说中的海上仙山。蓬莱仙境遥不可及，就像庾信永远回不去的江南故乡。}
\notetext{8}{提挈：带领、扶老携幼。挈（qiè），牵、扶。战乱中带着一家老小逃难，颠沛流离。}
\notetext{9}{关河：关塞与山河。年复一年地在关山之间奔波，写尽流亡的辛酸。}
\notetext{10}{契阔：离散、生死阻隔。出自《诗经》“死生契阔，与子成说”——生死离合，本不可追问天意。}
\notetext{11}{灵光岿然：汉代灵光殿在战火中独存，“岿然独存”成了一个成语。庾信自比灵光殿——旧人都凋零了，只有自己还孤独地立在这个世界上。}
\end{annotations}

\begin{appreciation}
《哀江南赋》是中国文学史上最沉痛的骈赋之一。骈赋，就是用两两相对的句子写成的长文，讲究对仗辞藻。以前骈文常被讥为"华丽的外壳"；庾信往这个外壳里装进了真痛。这段是全赋最沉痛的一段。

开篇八字定调："日暮途远，人间何世！"日暮途远出自伍子胥：被逼到绝路时他说"日暮途远，故倒行而逆施之"——太阳快落山，路看不到头。庾信借这四字说自己：人老了，家在万里之外，归途漫漫。"人间何世"是一句质问：这人间变成了什么世道？还记得《天问》吗？屈原对着天问了一百多个问题；庾信只问一句，里面全是答案。

"将军一去，大树飘零；壮士不还，寒风萧瑟"用了两个典故。大树将军是东汉的冯异，诸将争功时他独自坐在大树下；壮士是荆轲，唱着"风萧萧兮易水寒，壮士一去兮不复还"渡过易水。两个典故说的是同一件事：一去不还。南朝的将才壮士，都回不来了。

"舟楫路穷"以下四句写归路。传说有人乘木筏顺海漂到银河；海上还有蓬莱。庾信却说，银河不是乘筏能上的，蓬莱没有到达之期。回江南的路，比登天还难。最扎心的是他还在算——路都绝了，人还在想。

"提挈老幼，关河累年"写他自己的逃难：扶老携幼，在关山之间流亡一年又一年。"死生契阔"出自《诗经》，本是生死离合的誓言，如今应验了，却"不可问天"。末句用汉代灵光殿的典故：大乱后别的宫殿都毁了，只有灵光殿独存。庾信自比灵光殿：旧人零落殆尽，只有自己还孤零零地立着。活着，有时比死了更苦。

两百多年后杜甫写他："庾信平生最萧瑟，暮年诗赋动江关。"杜甫经历过国破，他懂。一个人的痛隔着两百年被另一个诗人接住，这就是文学。

妹妹，庾信是回不去，你是回得来。以后不管走多远，想家了随时回来，家里的饭桌永远有你的位置。
\end{appreciation}
